\section{Interpretation and reproducibility}\label{sec:limits}
\subsection{What an access-conditioned benchmark can establish}
The same ten-score threshold defines different experiments depending on whether its warm-up is scored. A common evaluation mask addresses that ambiguity directly. The same set of participant--model outcomes also gives different counts according to whether thresholding for a large decrease precedes or follows model averaging. Retaining participant, model and session keys until the final aggregation makes these choices reviewable. This matters when a benchmark serves both retrospective comparison and a decision about one deployed checkpoint.

Adaptive decision boundaries have a precedent in motor-imagery BCI. PMean updates an LDA bias from a running pooled feature mean under a fixed separation model, and its feedback experiment supplied feedback using the PMean rule itself~\cite{vidaurre2011}. The present median rule uses recorded deep-model scores and leaves all feature weights fixed. Here the participant cannot respond to a changed decision. Likewise, methods such as Average Thresholded Confidence seek to predict target accuracy in the absence of target labels~\cite{garg2022atc}. ACDD uses labels retrospectively for evaluation and diagnostic oracles. A fitted rule can improve without providing an unlabeled estimator of its own performance.

The scope is six public cohorts of cued binary left/right imagery, two backbones, one common sensor montage and a declared balanced prior. A different prevalence, a different cost for errors, class abstention or a new montage can change which comparison is relevant. Source validation was reused for checkpoint selection and calibration. Decreases in BA describe classification outcomes. They do not measure an injury or adverse clinical event. Prospective feedback, asynchronous or multiclass control, future users, and other time-series domains require separate experiments.

\subsection{Artifact, provenance and reproducibility}
The local artifact identifies the 30 paired-logit target files, labeled source-validation inputs and checkpoint hashes. With recorded scores, its first group of calculations reconstructs the protected decisions, common-suffix decisions and cross-fitted references. A source-only refit of Platt-v2 reproduces the persisted parameter values at full precision. Loading those verified snapshots regenerates the decision hashes, 59,520 candidate participant--model--policy rows and 1,860 paired utility rows. The source-only, EM and BBSE controls in Table~\ref{tab:policies} use this version. The protected quantile uses no Platt parameters.

The historical Platt coefficients were unavailable. Against archived BA, a legacy refit changed 116 of 14,880 protected participant--model--policy rows (62 source Platt and 54 EM, one EEGNet checkpoint). The largest absolute difference was 2.98 points. Separately, Platt-v2 changed 1,516 of 59,520 rows across four windows at tolerance $10^{-12}$, with a maximum of 6.47 points. In Table~\ref{tab:policies}, its protected EEGNet EM mean moved from $-2.512$ to $-2.484$ points and its count below $-2$ fell from 32 to 29. For Conformer source Platt, 19 participants had negative mean gain, compared with 22 historically. All eight gain intervals retained their direction. The median policy, its fixed-model counts and the raw-logit oracles use no fitted Platt map.

The archived discovery input begins at participant metrics. Dreyer starts from trial predictions, and the synthetic study regenerates scores from its design. Software, checkpoint locations and official acquisition paths are identified for the EEG-to-logit path. Raw-signal replay and end-to-end temporal invariance were not run. The historical registry for Fig.~\ref{fig:cohorts} names one CSV for all four Stieger rows but carries the digest of the earlier confirmation JSON. Both files are available. Their four means agree, while all four bootstrap intervals differ. The confirmation analysis resampled participant means and used order-statistic limits. Its protected gain intervals were [1.02, 1.81] points for EEGNet and [1.28, 2.25] for the Conformer. The plotted intervals use interpolated percentiles from the later participant bootstrap after averaging seeds and source conditions. Replaying the two score rules and that bootstrap reproduced all four plotted triples. The candidate manifest pairs each file with its own digest and records the bootstrap used for the plotted intervals. The author archive contains the verified frozen inputs. The reviewer package lists each restricted input by digest. Its distribution route and public artifact URL await author approval.

\subsection{Ethics}
The study reuses public de-identified recordings and involved no new recruitment. Original sources document oversight and consent for Cho2017, Lee2019-MI, Dreyer2023 and Stieger2021. The records reviewed for BNCI2014-001 and PhysioNetMI did not establish those details. The authors' institutional determination for secondary use awaits confirmation.

\section*{Acknowledgments and AI-generated content}
In earlier revisions, OpenAI Codex assisted with drafting and revising Sections I--VI, literature retrieval, code inspection and scripting for the frozen-score replay, the common-suffix calculation and table construction. It also assisted with numerical and layout checks and artifact documentation. The workflow and five result figures reuse the earlier manuscript's artwork. Edits to the workflow and participant figure clarify the score interface and the aggregation unit. In the preceding submission-preparation pass, its review of the scientific prose was limited to grammar. For this edition, Codex helped clarify trial eligibility and implement and test the reproduction package. Author review of the interpretation and final disclosure is pending.
